%==============================================================================
% PROGRESS REPORT - PROYEK PENGGANTI UAS
% Mata Kuliah Kecerdasan Buatan | Teknik Komputer - Universitas Indonesia
% Semester Genap 2025/2026
% Kelompok 06
%==============================================================================

\documentclass[11pt,a4paper,oneside]{report}

%------ Bahasa & Font --------------------------------------------------------
\usepackage{fontspec}
\setmainfont{Times New Roman}
\setmonofont{Consolas}[Scale=0.92]

%------ Geometri & Spasi -----------------------------------------------------
\usepackage[a4paper,margin=2.5cm]{geometry}
\usepackage{setspace}
\onehalfspacing
\setlength{\parskip}{0.5em}

%------ Paket Pendukung ------------------------------------------------------
\usepackage{graphicx}
\usepackage{amsmath,amssymb,amsfonts}
\usepackage{booktabs}
\usepackage{array}
\usepackage{tabularx}
\usepackage{longtable}
\usepackage{enumitem}
\usepackage{xcolor}
\usepackage{listings}
\usepackage{hyperref}
\usepackage{titlesec}
\usepackage{tcolorbox}
\usepackage{caption}

%------ Penyesuaian Bahasa Indonesia -----------------------------------------
\renewcommand{\chaptername}{BAB}
\renewcommand{\contentsname}{DAFTAR ISI}
\renewcommand{\listfigurename}{DAFTAR GAMBAR}
\renewcommand{\listtablename}{DAFTAR TABEL}
\renewcommand{\appendixname}{LAMPIRAN}
\renewcommand{\figurename}{Gambar}
\renewcommand{\tablename}{Tabel}

%------ Format Bab -----------------------------------------------------------
\titleformat{\chapter}[display]
  {\normalfont\Large\bfseries\centering}
  {\chaptername\ \thechapter}{12pt}{\Large}
\titlespacing*{\chapter}{0pt}{-20pt}{20pt}

%------ Hyperref -------------------------------------------------------------
\hypersetup{
  colorlinks=true,
  linkcolor=black,
  urlcolor=blue,
  citecolor=black,
  pdftitle={Progress Report - MK Kecerdasan Buatan}
}

%------ Listing untuk JSON ---------------------------------------------------
\lstdefinelanguage{json}{
  basicstyle=\ttfamily\footnotesize,
  showstringspaces=false,
  breaklines=true,
  frame=single,
  rulecolor=\color{gray!50},
  string=[s]{"}{"},
  stringstyle=\color{blue!60!black}
}

%------ Box Status -----------------------------------------------------------
\newcommand{\statusbox}[1]{\textbf{[STATUS: #1]}}

%==============================================================================
\begin{document}

%==============================================================================
% HALAMAN SAMPUL
%==============================================================================
\begin{titlepage}
\centering
\vspace*{1cm}

{\Large\bfseries PROGRESS REPORT \par}
{\large PROYEK PENGGANTI UAS \par}
\vspace{0.3cm}
{\large Mata Kuliah Kecerdasan Buatan \par}
{\large Semester Genap 2025/2026 \par}

\vspace{2cm}

\vspace{1.5cm}

{\LARGE\bfseries Klasifikasi Lalu Lintas Jaringan DDoS Menggunakan LSTM, GRU, dan Random Forest pada Dataset CIC-DDoS2019 \par}
\vspace{0.5cm}
{\large Perbandingan Arsitektur Temporal Deep Learning dan Machine Learning Tradisional untuk Deteksi Serangan Jaringan \par}

\vspace{2cm}

{\large Periode Laporan: \par}
{\large 28 April 2026 -- 19 Mei 2026 \par}

\vspace{1cm}

{\large Disusun oleh: \par}
\vspace{0.5cm}
\begin{tabular}{rl}
Calvin Wirathama Katoroy & NPM. 2306242395 \\
Syahmi Hamdani           & NPM. 2306220532 \\
Christover Angelo        & NPM. 2306220323 \\
Matthew Immanuel         & NPM. 2306221024 \\
\end{tabular}

\vspace{1.5cm}

{\large Dosen Pengampu: \par}
{\large Muhammad Firdaus Syawaludin Lubis, Ph.D. \par}

\vfill

{\large DEPARTEMEN TEKNIK KOMPUTER \par}
{\large FAKULTAS TEKNIK \par}
{\large UNIVERSITAS INDONESIA \par}
{\large MEI 2026 \par}
\end{titlepage}

%==============================================================================
% DAFTAR ISI
%==============================================================================
\tableofcontents

%==============================================================================
% RINGKASAN EKSEKUTIF
%==============================================================================
\chapter*{RINGKASAN EKSEKUTIF}
\addcontentsline{toc}{chapter}{Ringkasan Eksekutif}

Proyek ini bertujuan mengklasifikasikan lalu lintas jaringan sebagai normal atau serangan DDoS menggunakan pendekatan deep learning berbasis RNN. Tiga model dibandingkan: LSTM sebagai model utama, GRU sebagai baseline 1, dan Random Forest sebagai baseline 2. Pemilihan LSTM didasari sifat sekuensial lalu lintas jaringan---pola temporal antar paket merupakan sinyal kritis untuk membedakan serangan dari aktivitas benign.

Dataset yang digunakan adalah CIC-DDoS2019 versi \textit{cleaned} dari Kaggle (431.371 sampel, 17 jenis serangan DDoS, binary label: normal vs.\ ddos). Pipeline preprocessing mencakup: normalisasi label, seleksi 16 fitur CICFlowMeter, StandardScaler, stratified 70/15/15 split (seed 42), dan sliding window seq\_len=10 untuk model sekuensial.

Per laporan ini, seluruh training dan evaluasi telah selesai dilaksanakan. Ketiga model telah dilatih di Google Colab GPU dengan masing-masing minimal 6 eksperimen hyperparameter tuning dan dievaluasi pada test set. Hasil pada test set: LSTM (F1=99.76\%, FPR=0.46\%), GRU (F1=99.75\%, FPR=0.28\%), Random Forest (F1=99.92\%, FPR=0.15\%). Penulisan laporan akhir sedang berlangsung sebelum \textit{deadline} 25 Mei 2026.

\vspace{0.5em}
\noindent\textbf{Persentase progres keseluruhan:} \fbox{85\%}

%==============================================================================
% BAB I - PENDAHULUAN & DEFINISI MASALAH
%==============================================================================
\chapter{PENDAHULUAN \& DEFINISI MASALAH}
\label{bab:pendahuluan}

\section{Topik yang Dipilih}

Topik yang dipilih adalah \textbf{RNN / GRU / LSTM} (Topik 2). Lalu lintas jaringan pada dasarnya bersifat sekuensial---sebuah aliran jaringan terdiri dari serangkaian paket yang memiliki ketergantungan temporal satu sama lain. LSTM, sebagai varian Recurrent Neural Network (RNN) dengan mekanisme \textit{gating}, mampu menangkap pola temporal jangka panjang yang relevan dengan karakteristik serangan DDoS, seperti lonjakan volume paket yang terstruktur dalam waktu. Misalnya, sebuah lonjakan pada $t=47$ hanya bernilai anomali dalam konteks rentang $t=44..46$; pendekatan non-sekuensial seperti Random Forest tidak dapat menangkap konteks ini secara langsung.

Topik ini juga berkaitan langsung dengan penelitian laboratorium kami di Network Laboratory DTE FT UI (repo: \url{https://github.com/calvinkatoroy/wazuh-gan-ddos-research}), di mana LSTM berpotensi menggantikan Random Forest sebagai diskriminator dalam pipeline GANDD-Bridge pada SIEM Wazuh.

\section{Rumusan Masalah (Tentatif)}

\begin{enumerate}[label=RM\arabic*.,leftmargin=*]
    \item Seberapa baik model LSTM mampu mengklasifikasikan lalu lintas jaringan sebagai normal atau serangan DDoS pada dataset CIC-DDoS2019?
    \item Apakah LSTM memberikan performa yang lebih baik dibanding GRU dan Random Forest dalam hal accuracy, F1-score, dan False Positive Rate?
    \item Bagaimana pengaruh hyperparameter utama (hidden\_size, seq\_len, num\_layers) terhadap performa LSTM dalam konteks deteksi DDoS?
\end{enumerate}

\section{Ruang Lingkup \& Batasan Sementara}

Proyek ini \textbf{mencakup}: (1) binary classification lalu lintas jaringan (normal vs.\ ddos); (2) perbandingan tiga arsitektur model pada dataset benchmark CIC-DDoS2019; (3) hyperparameter tuning terdokumentasi (minimal 5 eksperimen per model) untuk klaim bonus B1; (4) analisis etika dan bias terkait keterbatasan dataset untuk klaim bonus B3.

Proyek ini \textbf{tidak mencakup}: (1) klasifikasi multi-kelas per jenis serangan DDoS; (2) deployment model ke sistem produksi atau integrasi langsung dengan Wazuh SIEM (cakupan penelitian lab terpisah); (3) pengumpulan dataset mandiri; (4) pengujian pada jenis serangan baru (QUIC-based DDoS, HTTP/2 RST flood) yang tidak ada dalam CIC-DDoS2019.

%==============================================================================
% BAB II - PENDEKATAN AI YANG DIRENCANAKAN
%==============================================================================
\chapter{PENDEKATAN AI YANG DIRENCANAKAN}
\label{bab:pendekatan}

\section{Studi Literatur yang Telah Dilakukan}

\begin{enumerate}[leftmargin=*]
    \item Sharafaldin, I., Lashkari, A.H., Hakak, S., \& Ghorbani, A.A. (2019). ``Developing Realistic Distributed Denial of Service (DDoS) Attack Dataset and Taxonomy.'' \textit{IEEE ICCST}. --- Paper orisinal CIC-DDoS2019; menjelaskan metodologi penangkapan \textit{traffic}, distribusi kelas, dan fitur CICFlowMeter yang menjadi dasar dataset kami.

    \item Hochreiter, S., \& Schmidhuber, J. (1997). ``Long Short-Term Memory.'' \textit{Neural Computation, 9}(8), 1735--1780. --- Paper LSTM orisinal; dasar teori mekanisme \textit{forget/input/output gate} yang mengatasi vanishing gradient pada sekuensi panjang.

    \item Cho, K., van Merrienboer, B., Gulcehre, C., Bahdanau, D., Bougares, F., Schwenk, H., \& Bengio, Y. (2014). ``Learning Phrase Representations using RNN Encoder-Decoder for Statistical Machine Translation.'' \textit{EMNLP}. --- Memperkenalkan GRU sebagai penyederhanaan LSTM; dasar arsitektur baseline 1 kami.

    \item Breiman, L. (2001). ``Random Forests.'' \textit{Machine Learning, 45}(1), 5--32. --- Paper orisinal Random Forest; dasar implementasi baseline 2 dengan \textit{bagging} ensemble dan random feature subsets.

    \item Roopak, M., Tian, G.Y., \& Chambers, J. (2019). ``Deep Learning Models for Cyber Security in IoT Networks.'' \textit{IEEE ICASI}. --- Menunjukkan keunggulan LSTM atas ML tradisional dalam deteksi intrusi jaringan; mendukung pilihan LSTM sebagai model utama.
\end{enumerate}

\section{Arsitektur Model yang Dipilih (Sementara)}

Tiga arsitektur diimplementasikan untuk perbandingan komparatif:

\textbf{LSTM (Model Utama):} Menerima input berupa sekuens aliran jaringan dengan panjang window \texttt{seq\_len=10} dan 16 fitur CICFlowMeter. LSTM dengan \texttt{hidden\_size=128} dan 2 \textit{layers} (dropout=0.3) diikuti oleh \textit{classification head} Linear(128$\to$64, ReLU, Dropout) dan Linear(64$\to$2). Hanya \textit{output} pada \textit{timestep} terakhir yang digunakan untuk klasifikasi.

\textbf{GRU (Baseline 1):} Arsitektur identik dengan LSTM, hanya mengganti \texttt{nn.LSTM} dengan \texttt{nn.GRU}. GRU memiliki lebih sedikit parameter (tidak ada \textit{memory cell} terpisah), sehingga lebih cepat untuk dilatih namun berpotensi kurang ekspresif untuk pola temporal kompleks.

\textbf{Random Forest (Baseline 2):} Menggunakan \textit{flat feature vector} (tanpa \textit{windowing}) dengan \texttt{n\_estimators=100} dan \texttt{class\_weight=`balanced'} untuk menangani ketidakseimbangan kelas.

\begin{figure}[h!]
\centering
\fbox{\parbox[c][6cm][c]{0.7\textwidth}{\centering
\textbf{Arsitektur LSTM (Model Utama)} \\[0.8em]
\texttt{Input}: (batch $\times$ seq\_len=10 $\times$ 16 fitur) \\[0.3em]
$\downarrow$ \\[0.3em]
\texttt{nn.LSTM}(hidden=128, layers=2, dropout=0.3, batch\_first=True) \\[0.3em]
$\downarrow$ ambil output $h_{t_{\text{last}}}$: (batch $\times$ 128) \\[0.3em]
$\downarrow$ \\[0.3em]
\texttt{Linear}(128 $\to$ 64) + \texttt{ReLU} + \texttt{Dropout}(0.3) \\[0.3em]
$\downarrow$ \\[0.3em]
\texttt{Linear}(64 $\to$ 2) \\[0.3em]
$\downarrow$ \\[0.3em]
\texttt{Output}: logit [normal, ddos]
}}
\caption{Sketsa arsitektur LSTM yang direncanakan. GRU identik kecuali mengganti \texttt{nn.LSTM} dengan \texttt{nn.GRU}.}
\label{fig:arsitektur-sementara}
\end{figure}

\section{Justifikasi Awal Pemilihan Pendekatan}

LSTM dipilih sebagai model utama karena dua alasan utama. Pertama, lalu lintas jaringan memiliki dimensi temporal yang kuat: sebuah serangan DDoS memanifestasikan dirinya sebagai lonjakan volume paket dalam rentang waktu, bukan sebagai satu paket tunggal yang anomali. LSTM dapat menangkap dependensi ini melalui mekanisme \textit{gating}-nya, sedangkan Random Forest memproses setiap sampel sebagai observasi independen. Kedua, proyek ini merupakan ekstensi dari penelitian laboratorium kami yang menggunakan Random Forest sebagai diskriminator dalam GANDD-Bridge; membandingkan LSTM/GRU dengan RF memberikan kontribusi langsung terhadap pertanyaan apakah arsitektur sekuensial dapat meningkatkan performa sistem deteksi tersebut.

GRU dipilih sebagai baseline pertama karena berbagi prinsip yang sama dengan LSTM (gated recurrent unit) namun dengan arsitektur yang lebih sederhana, memungkinkan perbandingan efisiensi-performa yang bermakna. Random Forest dipilih sebagai baseline kedua karena merupakan metode yang sudah terbukti (\textit{established}) dalam domain deteksi intrusi jaringan dan merepresentasikan pendekatan ML tradisional non-sekuensial yang digunakan pada sistem produksi laboratorium kami.

%==============================================================================
% BAB III - PROGRES PEKERJAAN SAAT INI
%==============================================================================
\chapter{PROGRES PEKERJAAN SAAT INI}
\label{bab:progres}

\section{Ringkasan Status per Tahap}

\begin{table}[h!]
\centering
\caption{Status per tahapan pekerjaan}
\label{tab:status-tahapan}
\begin{tabularx}{\textwidth}{lXl}
\toprule
\textbf{Tahap} & \textbf{Aktivitas} & \textbf{Status} \\
\midrule
Pencarian dataset      & Download CIC-DDoS2019 cleaned dari Kaggle (format Parquet, 431K+ sampel)       & [$\checkmark$] \\
EDA \& pra-pemrosesan  & Seleksi 16 fitur, binary label mapping, StandardScaler, stratified 70/15/15 split, sliding window seq\_len=10       & [$\checkmark$] \\
Implementasi model     & LSTM, GRU, RF modules di \texttt{src/models/}; training loop dan evaluasi selesai ditulis       & [$\checkmark$] \\
Training awal          & Selesai untuk ketiga model di Google Colab GPU       & [$\checkmark$] \\
Evaluasi \& metrik     & Confusion matrix, classification report, ROC/AUC per model telah digenerate       & [$\checkmark$] \\
Hyperparameter tuning  & 6 eksperimen per model (LSTM \& GRU), 5 eksperimen RF, terdokumentasi di CSV       & [$\checkmark$] \\
Error analysis         & Per-class error analysis via classification report dan confusion matrix       & [$\checkmark$] \\
Penulisan laporan akhir & Draft Bab I--IV selesai, Bab V (kesimpulan) sedang disusun       & [$\bullet$] \\
\bottomrule
\end{tabularx}
\end{table}

\noindent Keterangan: $\checkmark$ = Selesai \quad $\bullet$ = Sedang berjalan \quad $\square$ = Belum mulai.

\section{Dataset \& Pra-pemrosesan}

\statusbox{selesai}

Dataset CIC-DDoS2019 versi \textit{cleaned} dari Kaggle (\url{https://www.kaggle.com/datasets/dhoogla/cicddos2019}) telah berhasil diperoleh dalam format Parquet. Dataset ini merupakan versi yang sudah dibersihkan dari masalah kualitas data pada CSV orisinal CIC-UNB, yaitu nilai \texttt{Inf} yang dihasilkan CICFlowMeter pada aliran \textit{zero-duration} (pembagian dengan nol), nilai null, duplikat, dan kolom identifier (IP, port, timestamp) yang berpotensi menyebabkan \textit{data leakage}.

\textbf{Statistik dataset:}
\begin{itemize}[leftmargin=*]
    \item Total sampel setelah penggabungan: 431.371 baris
    \item Label setelah binary mapping: \texttt{normal} (Benign, $\approx$97.831 sampel) dan \texttt{ddos} (semua jenis serangan, $\approx$333.540 sampel)
    \item Ketidakseimbangan kelas ditangani dengan \texttt{class\_weight=`balanced'} untuk RF dan weighted loss untuk LSTM/GRU
\end{itemize}

\textbf{Pipeline pra-pemrosesan} (diimplementasikan dalam \texttt{src/preprocessing.py}):
\begin{enumerate}[leftmargin=*]
    \item Load semua file Parquet dan normalisasi nama label ke binary (\texttt{normal}/\texttt{ddos})
    \item Seleksi 16 fitur CICFlowMeter yang relevan secara temporal
    \item Eliminasi baris dengan nilai NaN atau Inf
    \item Stratified 70/15/15 split (seed 42) untuk menjaga proporsi kelas di setiap subset
    \item StandardScaler di-\textit{fit} pada \textit{train set} saja, ditransformasi ke val dan test (mencegah \textit{data leakage})
    \item Sliding window (seq\_len=10) untuk LSTM/GRU: menghasilkan tensor berukuran $(n, 10, 16)$
    \item \textit{Flat feature vector} untuk Random Forest: menghasilkan array $(n, 16)$
\end{enumerate}

\section{Implementasi Model}

\statusbox{selesai}

Seluruh modul kode telah ditulis dan siap dijalankan di Google Colab. Struktur direktori \texttt{src/} adalah sebagai berikut:

\begin{itemize}[leftmargin=*]
    \item \texttt{src/models/lstm\_model.py}: \texttt{LSTMClassifier} (\texttt{nn.Module} PyTorch)
    \item \texttt{src/models/gru\_model.py}: \texttt{GRUClassifier} (identik dengan LSTM, menggunakan \texttt{nn.GRU})
    \item \texttt{src/models/rf\_model.py}: \textit{wrapper} \texttt{RandomForestClassifier} dari scikit-learn
    \item \texttt{src/train.py}: training loop generik dengan logging per \textit{epoch} ke DataFrame
    \item \texttt{src/evaluate.py}: \texttt{classification\_report}, confusion matrix, ROC/AUC
    \item \texttt{src/preprocessing.py}: pipeline preprocessing lengkap (lihat Bab III.2)
    \item \texttt{src/utils.py}: seed setting, tqdm progress bar, plot helpers
    \item \texttt{config.yaml}: semua hyperparameter dan path (tidak ada \textit{magic number} dalam kode)
\end{itemize}

Notebook Colab tersedia di direktori \texttt{notebooks/}: \texttt{03\_lstm\_training.ipynb}, \texttt{04\_gru\_training.ipynb}, \texttt{05\_rf\_training.ipynb}, dan \texttt{08\_evaluation.ipynb}. Setiap notebook memanggil fungsi dari \texttt{src/} via \texttt{sys.path.insert}.

Repository: \url{https://github.com/calvinkatoroy/tugas-akhir-ai-kel-06}

\section{Hasil Eksperimen Sementara (Jika Ada)}

\statusbox{selesai}

Seluruh training telah selesai. Berikut hasil terbaik masing-masing model pada test set (stratified 15\% dari 431.371 sampel):

\begin{table}[h!]
\centering
\caption{Hasil evaluasi pada test set (\textit{preliminary} --- final)}
\label{tab:hasil-sementara}
\begin{tabular}{lccccc}
\toprule
\textbf{Model} & \textbf{Best Val F1} & \textbf{Test F1} & \textbf{Accuracy} & \textbf{FPR} & \textbf{AUC} \\
\midrule
LSTM (hidden=128, layers=2, seq\_len=10) & 0.9976 & 0.9976 & 99.63\% & 0.46\% & 99.96\% \\
GRU (hidden=256, layers=2, seq\_len=10)  & 0.9976 & 0.9975 & 99.61\% & 0.28\% & 99.96\% \\
Random Forest (n\_estimators=100)        & --     & 0.9992 & 99.87\% & 0.15\% & 99.98\% \\
\bottomrule
\end{tabular}
\end{table}

Temuan menarik: Random Forest mengungguli kedua model deep learning pada semua metrik. Hal ini akan dibahas dalam laporan akhir---RF kemungkinan diuntungkan karena fitur CICFlowMeter sudah merupakan \textit{engineered features} yang informatif, sehingga kemampuan LSTM menangkap dependensi temporal tidak memberikan keunggulan signifikan. Hyperparameter tuning didokumentasikan dalam \texttt{results/metrics\_summary.csv} (6 eksperimen LSTM, 6 eksperimen GRU, 5 eksperimen RF).

%==============================================================================
% BAB IV - PEMBAGIAN TUGAS & DEKLARASI KONTRIBUSI INDIVIDUAL
%==============================================================================
\chapter{PEMBAGIAN TUGAS \& DEKLARASI KONTRIBUSI INDIVIDUAL}
\label{bab:kontribusi}

\section{Calvin Wirathama Katoroy -- NPM. 2306242395}

\subsection*{IV.1.1 Komponen yang Dikerjakan}
Daftar konkret bagian proyek yang menjadi tanggung jawab penuh saya:
\begin{itemize}[leftmargin=*]
    \item \textbf{Kode/modul:} \texttt{src/models/lstm\_model.py}, \texttt{src/models/gru\_model.py}, \texttt{src/train.py}, \texttt{config.yaml}, \texttt{train\_local.py}
    \item \textbf{Bagian arsitektur:} Desain LSTM/GRU classifier (pilihan \texttt{hidden\_size}, \texttt{num\_layers}, \texttt{dropout}, classification head), training loop dengan logging per epoch, weighted loss untuk class imbalance
    \item \textbf{Bagian analisis:} Justifikasi pemilihan LSTM vs.\ GRU, rancangan hyperparameter grid, koneksi dengan penelitian GANDD-Bridge
    \item \textbf{Bagian laporan:} Bab I, Bab II.2, Bab II.3, Bab III.3, koordinasi keseluruhan dokumen
\end{itemize}

\subsection*{IV.1.2 Pernyataan Penggunaan AI \textit{Assistant}}
\begin{itemize}[leftmargin=*]
    \item \textbf{Tool yang dipakai:} Claude Sonnet 4.6 (Claude Code CLI, Anthropic)
    \item \textbf{Untuk bagian apa:} Debugging sesekali pada training loop dan path resolution di Google Colab
    \item \textbf{Pernyataan:} Saya memahami dan dapat mempertanggungjawabkan setiap baris kode dan teks yang saya klaim sebagai bagian saya.
\end{itemize}

%-------------------- Anggota 2 ----------------------
\section{Syahmi Hamdani -- NPM. 2306220532}

\subsection*{IV.2.1 Komponen yang Dikerjakan}
Daftar konkret bagian proyek yang menjadi tanggung jawab penuh saya:
\begin{itemize}[leftmargin=*]
    \item \textbf{Kode/modul:} \texttt{src/preprocessing.py}, \texttt{verify\_dataset.py}, \texttt{notebooks/01\_eda.ipynb}, \texttt{notebooks/02\_preprocessing.ipynb}
    \item \textbf{Bagian arsitektur:} Pipeline pra-pemrosesan data (StandardScaler fitted pada train set untuk mencegah data leakage, pembagian stratified 70/15/15 split, pembuatan sliding window sequence berukuran $(n, 10, 16)$), skrip validasi integritas dataset splits
    \item \textbf{Bagian analisis:} Analisis eksplorasi sebaran kelas (class imbalance), analisis korelasi fitur jaringan (Pearson correlation heatmap), penentuan 16 fitur temporal kritis dari CICFlowMeter
    \item \textbf{Bagian laporan:} Bab III.1 (Status per Tahap), Bab III.2 (Dataset \& Pra-pemrosesan), Lampiran A (Metadata JSON)
\end{itemize}

\subsection*{IV.2.2 Pernyataan Penggunaan AI \textit{Assistant}}
\begin{itemize}[leftmargin=*]
    \item \textbf{Tool yang dipakai:} Claude Sonnet 4.6 (Claude Code CLI, Anthropic)
    \item \textbf{Untuk bagian apa:} Debugging sesekali pada preprocessing pipeline dan validasi skrip dataset splits offline
    \item \textbf{Pernyataan:} Saya memahami dan dapat mempertanggungjawabkan setiap baris kode dan teks yang saya klaim sebagai bagian saya.
\end{itemize}

%-------------------- Anggota 3 ---------------------------------------------
\section{Christover Angelo -- NPM. 2306220323}

\subsection*{IV.3.1 Komponen yang Dikerjakan}
Daftar konkret bagian proyek yang menjadi tanggung jawab penuh saya:
\begin{itemize}[leftmargin=*]
    \item \textbf{Kode/modul:} \texttt{notebooks/09\_error\_analysis.ipynb}
    
    \item \textbf{Bagian arsitektur/eksekusi:} Menambahkan notebook analisis lanjutan untuk mengevaluasi performa model klasifikasi DDoS yang telah dilatih sebelumnya tanpa mengubah pipeline utama proyek maupun konfigurasi training model.
    
    \item \textbf{Bagian analisis:} Melakukan \textit{error analysis} terhadap model LSTM, GRU, dan Random Forest, termasuk visualisasi perbandingan akurasi, perhitungan \textit{error rate}, serta analisis \textit{False Positive} dan \textit{False Negative} pada sistem deteksi serangan DDoS.
    
    \item \textbf{Bagian laporan:} Analisis Error Model, Interpretasi Hasil Evaluasi Model, dan Pembahasan Kelebihan serta Kekurangan Tiap Model.
\end{itemize}

\subsection*{IV.3.2 Pernyataan Penggunaan AI \textit{Assistant}}
\begin{itemize}[leftmargin=*]
    \item \textbf{Tool yang dipakai:} Cursor
    
    \item \textbf{Untuk bagian apa:} Membantu Debugging dan penyusunan struktur analisis.
    
    \item \textbf{Pernyataan:} Saya memahami kode, visualisasi, dan isi analisis yang saya tambahkan sebagai bagian kontribusi individu saya pada proyek ini.
\end{itemize}

%-------------------- Anggota 4 ------
\section{Matthew Immanuel -- NPM. 2306221024}

\subsection*{IV.4.1 Komponen yang Dikerjakan}
Daftar konkret bagian proyek yang menjadi tanggung jawab penuh saya:
\begin{itemize}[leftmargin=*]
    \item \textbf{Kode/modul:} \texttt{notebooks/05\_rf\_training.ipynb}, \texttt{notebooks/08\_evaluation.ipynb}, manajemen pencatatan pada \texttt{metrics\_summary.csv}, dan direktori \texttt{results/figures/}
    \item \textbf{Bagian arsitektur/eksekusi:} Menjalankan eksperimen \textit{tuning hyperparameter} model \textit{baseline} Random Forest sebanyak 5 kali iterasi (variasi \texttt{n\_estimators} dan \texttt{max\_depth}) dan mengelola pemuatan (\textit{loading}) \textit{checkpoint} model untuk evaluasi.
    \item \textbf{Bagian analisis:} Mengeksekusi komparasi performa final ketiga arsitektur model (LSTM, GRU, Random Forest), menghasilkan tabel perbandingan, serta membuat visualisasi \textit{Confusion Matrix} dan \textit{ROC Curve overlay}.
    \item \textbf{Bagian laporan:} Desain Eksperimen Baseline, Hasil Evaluasi dan Analisis Komparatif, Kesimpulan Perbandingan Model
\end{itemize}

\subsection*{IV.4.2 Pernyataan Penggunaan AI \textit{Assistant}}
\begin{itemize}[leftmargin=*]
    \item \textbf{Tool yang dipakai:} Gemini, ChatGPT
    \item \textbf{Untuk bagian apa:} \textit{Debugging} sintaks Python, perbaikan \textit{error} saat memuat \textit{checkpoint} antar-\textit{environment}, dan penyesuaian parameter visualisasi grafik Matplotlib.
    \item \textbf{Pernyataan:} Saya memahami kode dan teks yang saya klaim sebagai bagian saya.
\end{itemize}

%==============================================================================
% LAMPIRAN
%==============================================================================
\appendix
\renewcommand{\thechapter}{\Alph{chapter}}
\titleformat{\chapter}[display]
  {\normalfont\Large\bfseries\centering}
  {\appendixname\ \thechapter}{12pt}{\Large}

%------ Lampiran A: Metadata JSON (versi sementara) -------------------------
\chapter{Berkas Metadata Proyek (Versi Sementara)}
\label{lamp:metadata}

Berikut adalah isi berkas \texttt{metadata-kel06.json} versi saat ini.

\begin{lstlisting}[language=json]
{
  "kode_kelompok": "K06",
  "topik_pilihan": 2,
  "nama_topik": "RNN / GRU / LSTM",
  "judul_proyek": "Klasifikasi Lalu Lintas Jaringan DDoS Menggunakan LSTM, GRU, dan Random Forest pada Dataset CIC-DDoS2019",
  "metode_ai": [
    {"nama": "LSTM",          "peran": "utama"},
    {"nama": "GRU",           "peran": "baseline"},
    {"nama": "Random Forest", "peran": "baseline"}
  ],
  "datasets": [
    {
      "nama": "CIC-DDoS2019 (Kaggle cleaned)",
      "tipe": "publik",
      "url_sumber_publik": "https://www.kaggle.com/datasets/dhoogla/cicddos2019",
      "jumlah_sampel": 431371,
      "jumlah_kelas": 2
    }
  ],
  "metrik_utama": ["Accuracy", "Precision", "Recall", "F1-Score", "FPR", "AUC-ROC"],
  "klaim_bonus": ["B1", "B3"],
  "tautan": {
    "github": "https://github.com/calvinkatoroy/tugas-akhir-ai-kel-06",
    "colab": "https://colab.research.google.com/drive/1rhoLvdyzKTdK9sGd4dRjkglO-lo2bBHB"
  },
  "ai_assistant_usage": {
    "digunakan": true,
    "tools": ["Claude Sonnet 4.6 (Claude Code CLI, Anthropic)"],
    "estimasi_persentase": 10
  }
}
\end{lstlisting}

%------ Lampiran B: Tautan Repositori ---------------------------------------
\chapter{Tautan Repositori \& Sumber}
\label{lamp:tautan}

\begin{description}[leftmargin=5cm,labelwidth=4.5cm]
    \item[GitHub repo:] \url{https://github.com/calvinkatoroy/tugas-akhir-ai-kel-06}
    \item[Google Colab:] \url{https://colab.research.google.com/drive/1rhoLvdyzKTdK9sGd4dRjkglO-lo2bBHB}
    \item[Dataset (Kaggle):] \url{https://www.kaggle.com/datasets/dhoogla/cicddos2019}
    \item[Dataset asli (CIC):] \url{https://www.unb.ca/cic/datasets/ddos-2019.html}
    \item[Lab research repo:] \url{https://github.com/calvinkatoroy/wazuh-gan-ddos-research}
\end{description}

\end{document}
